% Summary for the tissium-automation profile: AI & Automation Lead, TISSIUM
% (Paris, hybrid).
%
% Written to carry what the bullets below cannot, so it repeats none of
% them:
%   - what Fraunhofer IPA is. A French reader in a medtech company has no
%     reason to know that "IPA" is the Fraunhofer Institute for
%     Manufacturing Engineering and Automation (its official English name;
%     German: Fraunhofer-Institut für Produktionstechnik und
%     Automatisierung). That name is the posting's "manufacturing contexts"
%     and its "automation", on the first line.
%   - "from the idea to adoption": the posting's "full lifecycle:
%     identification through adoption". The coding-agent bullet is the
%     evidence -- started unasked, now a production service its researchers
%     use -- so the summary names the span and leaves the story to it.
%   - the teams: quantum computing, biotech, mechanical engineering and
%     electroplating (owner, 2026-09-23) -- "stakeholders across functions".
%   - agent skills for the quantum computing team (owner, 2026-09-23).
%   - Python in every project (owner, 2026-09-23).
%   - four years of teaching (section_teaching_mentoring.tex; that section
%     is not input here, this clause stands in for it) -- "train teams".
%   - the owner's public repositories, both checkable through the header's
%     GitHub link: ANTfrastructure's planner/executor loop
%     (docs/windows-agentic-loop.md) and DocumANTation, which typesets this
%     page -- the nearest thing on it to "document-heavy workflows".
% Earlier versions of this paragraph go below as comments, never deleted.
%
% (-1) 2026-09-28, first version: six lines in English, "this CV." alone on
% the last. The Ollama/Open WebUI stack went: this posting asks for no
% self-hosting, and the clause was exactly the overrun. In German
% "Quantencomputing-Team" also stood 27pt into the margin -- a word with an
% explicit hyphen breaks only there -- so the German now ends the clause on
% "Quantencomputing", which TeX can hyphenate. Kept:
%   EN: At Fraunhofer IPA, the Fraunhofer Institute for Manufacturing Engineering and Automation, I replace repetitive work with AI tools and automations that teams use, and carry them from the idea to adoption. I work with teams in quantum computing, biotech, mechanical engineering and electroplating, build agent skills for our quantum computing team, use Python in every project, and taught computer organization to undergraduates for four years. Publicly, I maintain a planner/executor loop for autonomous agents, an Ollama and Open WebUI stack for Nextcloud and the documentation pipeline that typeset this CV.
%   DE: Am Fraunhofer IPA, dem Fraunhofer-Institut für Produktionstechnik und Automatisierung, ersetze ich repetitive Arbeit durch KI-Werkzeuge und Automatisierungen, die Teams nutzen, und begleite sie von der Idee bis zur Einführung. Ich arbeite mit Teams aus Quantencomputing, Biotechnologie, Maschinenbau und Galvanik, baue Agent-Skills für unser Quantencomputing-Team, nutze Python in jedem Projekt und habe vier Jahre Rechnerorganisation gelehrt. Öffentlich pflege ich eine Planner/Executor-Schleife für autonome Agenten, einen Ollama- und Open-WebUI-Stack für Nextcloud und die Dokumentationspipeline, die diesen Lebenslauf gesetzt hat.
\IfLanguageName{english}{
    \par{At Fraunhofer IPA, the Fraunhofer Institute for Manufacturing Engineering and Automation, I replace repetitive work with AI tools and automations that teams use, and carry them from the idea to adoption. I work with teams in quantum computing, biotech, mechanical engineering and electroplating, build agent skills for our quantum computing team, use Python in every project, and taught computer organization to undergraduates for four years. Publicly, I maintain a planner/executor loop for autonomous agents and the documentation pipeline that typeset this CV.}
}{
    \par{Am Fraunhofer IPA, dem Fraunhofer-Institut für Produktionstechnik und Automatisierung, ersetze ich repetitive Arbeit durch KI-Werkzeuge und Automatisierungen, die Teams nutzen, und begleite sie von der Idee bis zur Einführung. Ich arbeite mit Teams aus Quantencomputing, Biotechnologie, Maschinenbau und Galvanik, baue Agent-Skills für unser Team im Quantencomputing, nutze Python in jedem Projekt und habe vier Jahre Rechnerorganisation gelehrt. Öffentlich pflege ich eine Planner/Executor-Schleife für autonome Agenten und die Dokumentationspipeline, die diesen Lebenslauf gesetzt hat.}
}
